Given a sequence $a_n$ of real or complex numbers, we can consider the sequence
of the so-called \emph{partial sums}%
\mymargin{partial sums}%
\index{sums!partial}
\[
  S_n = \sum_{k=0}^{n-1} a_k.
\]
\mynote{The sum starts from $k=0$ because we want to use all the terms of the sequence, but in many cases, the sequence will have its first index at $k=1$, and thus it will be natural to start from $k=1$. The last term is $n-1$ so that $S_n$ is the sum of the first $n$ terms of the sequence. Nothing would change if we chose $n$ instead of $n-1$.}
We can write more concisely $S_n = \sum a_n$.
Intuitively, it means summing the terms of the sequence $a_k$ for $k$ starting
from $0$ up to $k=n-1$:
\[
  \sum_{k=0}^{n-1} a_k = a_0 + a_1 + a_2 + \dots + a_{n-1}.
\]
Formally, the sum $S_n=\displaystyle \sum_{k=0}^n a_k$ is defined recursively by
the following relations:
\[
  \begin{cases}
    S_0 = 0, \\
    S_{n+1} = S_n + a_{n}.
  \end{cases}
\]

The numbers $a_n$ are called \emph{terms}%
\mymargin{terms}%
\index{term} of the series.
\index{terms!of a series}\index{series!terms}%
If the sequence of partial sums admits a limit, the limit is called the
\emph{sum}%
\mymargin{sum}%
\index{sum}
\index{sum!of a series}\index{series!sum}%
of the series and is denoted by
\[
  \sum_{k=0}^{+\infty} a_k 
  = \lim_{n\to +\infty} S_n 
  = \lim_{n\to+\infty} \sum_{k=0}^{n-1} a_k.
\]

The terminology already introduced for sequences also applies to series, which
are indeed sequences themselves. In particular, a series can be convergent,
divergent, or indeterminate. This is the \emph{nature of the series}.
\mymargin{nature}%
\index{nature!of a series}%
\index{series!nature}%

More generally, one can consider the sum starting from a certain index $m\in
\ZZ$. Given $m$, one can, for example, consider the series:
\[
  S_n = \sum_{k=m}^{n-1} a_k
\]
which is defined for every $n\in\NN$, $n\ge m$ as
\[
  S_n = a_m + a_{m+1} + \dots + a_{n-1}.
\]

\begin{example}
Consider the series $S_n$ defined as the sum of natural numbers from $1$ to $n$:
\[
  S_n 
  = \sum_{k=0}^{n-1} (k+1)
  = \sum_{k=1}^{n} k
  = 1 + 2 + \dots + n.
\]
We know that (exercise~\ref{ex:somma_lineare})
\[
  S_n = \frac{n(n+1)}{2}.
\]
And by the definition of limit, it can be verified that $S_n \to +\infty$.

This is expressed by saying that the series $\sum n$ is divergent:
\[
  \sum_{k=0}^{+\infty} (k+1)
  = \sum_{k=1}^{+\infty} k 
  = +\infty.
\]
\end{example}

\begin{example}[the geometric series]
Given $q \in \RR$ or $q\in \CC$, to the sequence
$  a_n = q^n$
of terms
\[
  a_0 = 1,\qquad
  a_1 = q,\qquad
  a_2 = q^2,\qquad
  a_3 = q^3, \dots
\]
is associated the
\emph{geometric series}%
\mymargin{geometric series}%
\index{series!geometric}
$\sum q^n$
whose partial sums are
\[
  S_0 = 0, \qquad
  S_1 = 1, \qquad
  S_2 = 1 + q, \qquad
  S_3 = 1 + q + q^2, \qquad
  \dots
\]
\end{example}

The following theorem shows us how, for different values of $q$, the geometric
series assumes all possible natures: convergent, divergent, indeterminate.

\begin{theorem}[sum of the geometric series]
\label{th:serie_geometrica}%
\mymark{***}%
\mymargin{sum of the geometric series}%
\index{sum!of the geometric series}%
Let $q\in \CC$. If $q\neq 1$, we have
\[
 \sum_{k=0}^{n-1} q^k  = \frac{1-q^n}{1-q}.
\]
If $\abs{q} < 1$, the geometric series converges:
\[
 \sum_{k=0}^{+\infty} q^k = \frac{1}{1-q}.
\]

If $\abs{q}\ge 1$, the series does not converge, but we must distinguish whether
the series is considered with real or complex values to determine its nature.

If we consider the series with real values (thus $q\in \RR$), if $q\ge 1$, the
series diverges to $+\infty$; if $q\le -1$, the series is indeterminate.

If we consider the series with complex values $q\in \CC$, if $\abs{q}>1$, the
series diverges to $\infty$, the same happens if $q=1$. If $\abs{q}=1$ but
$q\neq 1$, the series is indeterminate.
\end{theorem}
%
\begin{proof}
The first result concerns a finite sum. We have
\[
  (1-q)\cdot \sum_{k=0}^{n-1} q^k
  = \sum_{k=0}^{n-1} q^k - q \cdot \sum_{k=0}^{n-1} q^k
  = \sum_{k=0}^{n-1} q^k - \sum_{k=1}^{n} q^k
  = 1 - q^n
\]
from which, if $q\neq 1$, the desired result is obtained.

Taking the limit as $n\to +\infty$, if $\abs{q} < 1$, it is noted that
$\abs{q^n} = \abs{q}^{n} \to 0$, and the series thus converges to $\frac
1{1-q}$.

Now consider the series with real values. If $q>1$, we observe that $q^n\to
+\infty$, and thus the series diverges to $+\infty$ (in fact, in this case,
$1-q$ is negative). If $q=1$, we have $q^k=1$, and thus
$\displaystyle\sum_{k=0}^{n-1} q^k = n \to +\infty$.

For the other cases, we observe that
\[
  \sum_{k=0}^{n-1} q^k = \frac{1-q^n}{1-q}
  = \frac{1}{1-q} - \frac{1}{1-q}\cdot q^k
\]
and thus the nature of the series coincides with the nature of the sequence
$q^k$. If we consider the sequence with complex values and if $\abs{q}>1$, the
series is divergent to $\infty$ since $\abs{q^k}=\abs{q}^k\to +\infty$. If
instead, we consider the series with real values and $q<-1$ (the case $q\ge 1$
has already been considered), the series is indeterminate since in absolute
value it tends to $+\infty$, but the sign of the even terms is opposite to the
sign of the odd terms.

We still need to consider the case $\abs{q}=1$, $q\neq 1$. If $q$ is real, there
is only the case $q=-1$, and we know that $(-1)^k$ is indeterminate. Even if $q$
is complex, we want to show that the sequence $q^k$ is indeterminate: if it were
not, we would have $q^k\to \ell$ with $\abs{\ell}=1$ since $\abs{q^k}=1$, and
taking the limit in the equality
\[
  q^{k+2} = q\cdot q^{k+1}
\]
we would have $\ell = q\cdot \ell$, from which (dividing by $\ell$) we obtain
$q=1$. Thus, if $q\neq 1$, the sequence $q^k$ is indeterminate, and so is the
geometric series.
\end{proof}

\begin{theorem}[linearity of the sum]
\label{th:linearita_somma_serie}%
\mymargin{linearity of the infinite sum}%
\index{linearity of the infinite sum}
If $\sum a_n$ and $\sum b_n$ are convergent, then for every $\lambda,\mu\in \RR$
(or $\CC$), $\sum (\lambda a_n + \mu b_n)$ is also convergent, and we have
\[
 \sum_{k=0}^{+\infty} (\lambda a_n + \mu b_n)
 = \lambda \sum_{k=0}^{+\infty} a_n  + \mu \sum_{k=0}^{+\infty} b_n.
\]
\end{theorem}
%
\begin{proof}
If $S_n$ and $R_n$ are the partial sums of the series $\sum a_n$ and $\sum b_n$,
then the partial sums of the series $\sum (\lambda a_n + \mu b_n)$ are $\lambda
S_n + \mu R_n$ (since the distributive and commutative properties hold for
finite sums). But if $S_n \to S$ and $R_n \to R$, then $\lambda S_n + \mu R_n
\to \lambda S + \mu R$.
\end{proof}

We observe that series (like sequences) form a vector space in which the
operations of sum and scalar multiplication are performed term by term: $\sum
a_n + \sum b_n = \sum (a_n + b_n)$, $\lambda \sum a_n = \sum (\lambda a_n)$. The
previous theorem then tells us that convergent series (like convergent
sequences) form a subspace, and the sum of the series (like the limit of the
sequence) is a linear operator defined on this subspace.

\begin{theorem}[necessary condition for convergence]
\mymark{***}
\mymargin{necessary condition}%
\index{necessary condition}
If the series $\sum a_n$ converges, then $a_n \to 0$.
\end{theorem}
%
\begin{proof}
\mymark{***}
If the series $\sum a_n$ converges, it means that the partial sums $S_n =
\displaystyle\sum_{k=0}^{n-1} a_k$ converge: $S_n \to S$. But then
\[
  a_n = S_n - S_{n-1} \to S - S = 0.
\]
\end{proof}

\begin{theorem}[stability of the nature]%
\index{nature!of a series}%
\index{stability!of the nature}%
\index{series!stability of the nature}%
\index{series!that differ on a finite number of terms}%
If the two sequences $a_n$ and $b_n$ differ only on a finite number
\mymargin{stability of the nature}%
of terms, then the corresponding series $\sum a_n$ and $\sum b_n$ have the same
nature.
\end{theorem}
%
\begin{proof}
If the sequences differ on a finite number of terms, it means that there exists
an $N\in \NN$ such that for every $k>N$, we have $a_k=b_k$. Thus, if we denote
by $S_n = \sum_{k=0}^{n-1} a_k$ and $R_n = \sum_{k=0}^{n-1} b_k$ the
corresponding sequences of partial sums, for every $n>N$, we have
\[
  S_n - R_n
    = \sum_{k=0}^{n-1} a_k - \sum_{k=0}^{n-1} b_k
    = \sum_{k=0}^{N-1} (a_k - b_k) 
    = C
\]
where $C$ is a constant independent of $n$. Therefore, for $n>N$, we have
\[
  S_n = R_n + C.
\]
If the limit of $R_n$ does not exist, then the limit of $S_n$ does not exist
either (otherwise, since $R_n = S_n -C$, the limit of $R_n$ should also exist).
If the limit of $R_n$ is infinite, then the limit of $S_n$ is equal to the limit
of $R_n$. And if the limit of $R_n$ is finite, then the limit of $S_n$ is also
finite.

Thus, the nature of the sequence $S_n$ is the same as that of the sequence
$R_n$, i.e., the two series have the same nature.
\end{proof}

If a series has its first term with an index different from $0$, one can always
reduce it (with a change of variable) to a series whose index starts from zero.
For example (making the change of variable $j=k-1$ from which $j=0$ when $k=1$
and remembering that the index used in the sums of the series is a dummy
variable):
\[
 \sum_{k=1}^{+\infty} \frac{1}{2^k}
 = \sum_{j=0}^{+\infty} \frac{1}{2^{j+1}}
 = \sum_{k=0}^{+\infty} \frac{1}{2^{k+1}}.
\]

It should also be noted that according to the previous theorem, the first index
from which one starts summing is not relevant concerning the nature of the
series. However, if the series is convergent, its sum may vary, for example:
\[
 \sum_{k=1}^{+\infty} \frac{1}{2^k}
 = \enclose{\sum_{k=0}^{+\infty} \frac{1}{2^k}} - 2^0.
\]

Note well: in many books, $\infty$ is written instead of $+\infty$. It is
therefore very common to omit the $+$ sign in front of $\infty$ in the
terminology of series (and also sequences) since the indices are understood to
be natural numbers, and thus $-\infty$ would not make sense.

However, there are cases where it may be useful to use negative indices, for
example, if $a_k$ is defined for every $k\in \ZZ$, one could define (but we will
not do so):
\[
  \sum_{k=-\infty}^{+\infty} a_k
  = \sum_{k=0}^{+\infty} a_k +
  \sum_{k=1}^{+\infty} a_{(-k)}
\]
requiring that both series on the right side of the equality exist and do not
have infinite sums of opposite signs.

\begin{theorem}[tail of a convergent series]
\label{th:coda}%
\mymark{*}%
Let $\sum a_n$ be a convergent series. Then
\[
  \lim_{n\to +\infty} \sum_{k=n}^{+\infty} a_k = 0.
\]
\end{theorem}
%
\begin{proof}
\mymark{*}
Let
\[
  S_n = \sum_{k=0}^{n-1} a_k,
\]
by the definition of a convergent series, we know that there exists a finite $S$
such that $S_n \to S$. We then observe that
\[
  \sum_{k=n}^{+\infty} a_k = \lim_{N\to+\infty} \sum_{k=n}^{N-1} a_k
   = \lim_{N\to +\infty} S_N - S_n = S - S_n
\]
and, for $n\to +\infty$, we obviously have $S - S_n \to S - S = 0$.
\end{proof}